\clearpage\fancyhead[L]{\small GAUSS AN BODE.}\hypertarget{native-404-part-lower}{}\label{native-404-part-lower}
\begin{center}\large GAUSS AN BODE.\end{center}
\begin{center}\hrule\vspace{4pt}\fontsize{10}{12}\selectfont
Astronomisches Jahrbuch für 1820. Seite 201..205. Berlin 1817.\par\vspace{4pt}\hrule\vspace{4pt}\end{center}
\begin{flushright}Göttingen, den 17. März 1817.\end{flushright}
\par\indent Da die Theorie der Bewegung von Pallas hinreichend ausgearbeitet ist, habe ich Herrn Dr. \textsc{Tittel} aus Erlau, der sich anderthalb Jahre hier aufgehalten hat, veranlasst, die Ephemeride für 1818..1819 schon \SourceLeaf{405} jetzt zu berechnen. Ich freue mich, Ihnen eine Abschrift für Ihr nächstes Jahrbuch beizulegen.

\par\indent War der Himmel im vergangenen Winter bei Ihnen für Beobachtungen ebenso ungünstig wie hier? Nur höchst selten war die Sonne oder ein Stern zu sehen. Am ganzen 19. November des vergangenen Jahres schneite es, und von der Sonnenfinsternis war nichts zu sehen. Um die Mitte der Finsternis war eine gewisse Abnahme der Tageshelligkeit zu bemerken, doch war daran nichts Außergewöhnliches. Jedenfalls ist es seither auch ohne Sonnenfinsternis bei Schneegestöber öfter ebenso dunkel gewesen, selbst um Mittag. Bei der Mondfinsternis vom 4. Dezember war der Himmel etwas weniger ungünstig. Anfang und Ende blieben zwar unsichtbar, aber die folgenden Phasen habe ich, zum Teil zwischen rasch vorüberziehenden Wolken, ganz gut beobachtet.

\Needspace{15\baselineskip}\par\indent 4. Dezember 1816 auf der neuen Sternwarte.

\par\indent Eintritte.
\begin{center}\fontsize{10}{12}\selectfont\setlength{\tabcolsep}{3pt}\begin{tabular}{llrrrl}
Tycho 1. Rand & $\cdot\;\cdot\;\cdot\;\cdot\;\cdot\;\cdot\;\cdot\;\cdot\;\cdot\;\cdot\;\cdot\;\cdot$ &  $8^{\mathrm u}$ & $17^{\mathrm m}$ & $2^{\mathrm s}$ & W. Z.\\
Tycho 2. Rand & $\cdot\;\cdot\;\cdot\;\cdot\;\cdot\;\cdot\;\cdot\;\cdot\;\cdot\;\cdot\;\cdot\;\cdot$ &  & $18$ & $44$ & —\\
Kepler Mitte & $\cdot\;\cdot\;\cdot\;\cdot\;\cdot\;\cdot\;\cdot\;\cdot\;\cdot\;\cdot\;\cdot\;\cdot\;\cdot$ &  & $29$ & $24$ & —\\
Copernicus 1. Rand & $\cdot\;\cdot\;\cdot\;\cdot\;\cdot\;\cdot\;\cdot\;\cdot\;\cdot\;\cdot\;\cdot$ &  & $40$ & $28$ & —\\
Copernicus 2. Rand & $\cdot\;\cdot\;\cdot\;\cdot\;\cdot\;\cdot\;\cdot\;\cdot\;\cdot\;\cdot\;\cdot$ &  & $45$ & $11$ & —\\
Aristarch 1. Rand & $\cdot\;\cdot\;\cdot\;\cdot\;\cdot\;\cdot\;\cdot\;\cdot\;\cdot\;\cdot\;\cdot$ &  & $50$ & $28$ & —\\
Dionysius Mitte & $\cdot\;\cdot\;\cdot\;\cdot\;\cdot\;\cdot\;\cdot\;\cdot\;\cdot\;\cdot\;\cdot\;\cdot$ &  $9$ & $0$ & $46$ & —\\
Censorinus Mitte & $\cdot\;\cdot\;\cdot\;\cdot\;\cdot\;\cdot\;\cdot\;\cdot\;\cdot\;\cdot\;\cdot$ &  & $5$ & $44$ & —\\
\end{tabular}\end{center}

\Needspace{7\baselineskip}\begin{flushright}Göttingen, den 15. August 1817.\end{flushright}

\par\indent Der Bau der hiesigen neuen Sternwarte ist zwar noch nicht ganz abgeschlossen, aber doch so weit fortgeschritten, dass in zwei Zimmern beobachtet werden kann. Bis die neuen festen Meridianinstrumente eintreffen und aufgestellt sind, werden die Beobachtungen mit dem 12-zölligen \textsc{Reichenbach}schen Repetitionskreis die wichtigsten sein. Ich teile Ihnen die damit gewonnenen Ergebnisse der Beobachtungen des Polarsterns und der Sonne beim letzten Solstitium mit. Zunächst muss ich bemerken, dass nach meinen geodätischen Messungen die Mitte der neuen Sternwarte $6''1$ weiter südlich und um $1^{\mathrm s}90$ in Zeit weiter östlich liegt als die Mitte der alten Sternwarte. Ferner habe ich am Instrument selbst eine kleine Änderung vorgenommen. Bei meinen früheren Beobachtungen benutzte ich stets den Messingring, der auf der Objektivseite des Fernrohrs aufgesteckt wird und sowohl das Übergewicht des prismatischen Okulars als auch die Klammer ausgleicht. Man hatte vermutet, dass dieser Ring das Fernrohr durchbiegen könnte (oder genauer gesagt, dass er die Objektivseite stärker durchbiegen könnte als das Okular die andere Seite). Die Folge wäre, dass das Instrument alle Zenitdistanzen zu \emph{klein} angeben müsste. Man mutmaßte, dass dies vielleicht die Ursache des Unterschieds sei, den man bei der Bestimmung der Polhöhe je nach Beobachtung von Zirkumpolarsternen oder der Sonne findet. Wäre diese Vermutung begründet, so müsste der Unterschied verschwinden, wenn man das Gegengewicht wegnimmt. Vielmehr müsste sich der Unterschied dadurch umkehren, da jetzt eindeutig die Okularseite des Fernrohrs das Übergewicht hat. Aber meine neuen Beobachtungen, bei denen jenes Gewicht entfernt und stattdessen ein Gewicht am Alhidadenkreis selbst angeschraubt ist, sprechen ganz dagegen und scheinen mir eindeutig zu beweisen, dass die Durchbiegung des Fernrohrs überhaupt nicht wahrnehmbar ist. Zum einen ergeben die neuen Beobachtungen des Polarsterns und der Sonne wieder genau denselben Unterschied der \SourceLeaf{406}Polhöhe und in derselben Richtung wie die früheren. Zum anderen stimmt die absolute Bestimmung der Polhöhe der neuen Sternwarte genau mit der von der alten Sternwarte übertragenen überein.

\Needspace{15\baselineskip}\par\indent\emph{Beobachtungen des Polarsterns bei der unteren Kulmination.}

\begin{center}\fontsize{10}{12}\selectfont\setlength{\tabcolsep}{3pt}
\begin{tabular}{l|c|c|c}
1817 & Anzahl & \shortstack{Wahre Zenitdistanz\\im Meridian} & Polhöhe \\
\hline
März 28 & 16 & $40^\circ\ 8'\ 13''18$ & $51^\circ\ 31'\ 49''33$ \\
April\ \ 1 & 22 & $15.86$ & $47.84$ \\
\ \ \ \ \ \ \ \ 2 & 24 & $13.87$ & $50.13$ \\
\ \ \ \ \ \ \ \ 8 & 28 & $15.76$ & $50.03$ \\
\ \ \ \ \ \ \ 11 & 26 & $16.16$ & $50.51$ \\
\ \ \ \ \ \ \ 13 & 12 & $20.51$ & $46.75$ \\
Mai\ \ \ 7 & 30 & $23.90$ & $49.60$ \\
Mittel & 158 & $\cdot\ \cdot\ \cdot\ \cdot\ \cdot$ & $51^\circ\ 31'\ 49''42$ \\
\end{tabular}
\end{center}

\par\indent Für die Berechnung der Refraktion wurde \textsc{Carlini}s Tafel zugrunde gelegt, für die Deklination des Polarsterns die von \textsc{Bessel}. Der Beobachtungsort liegt $0''12$ nördlich der Mitte. Die Polhöhe in der Mitte beträgt also nach den obigen Beobachtungen $51^\circ\ 31'\ 49''3$.

\par\indent Meine Beobachtungen des Polarsterns bei der unteren Kulmination von 1813 ergeben, mit denselben Elementen reduziert, für die alte Sternwarte die Polhöhe $51^\circ\ 31'\ 55''48$. Die Übertragung liefert also dasselbe Ergebnis wie die unmittelbaren Beobachtungen.

\par\indent Die Witterung war in diesem Jahr für die Beobachtungen der Sommersonnenwende günstiger als je zuvor. Der Beobachtungsort liegt $0''16$ südlich der Mitte.

\Needspace{23\baselineskip}\par\indent\emph{Beobachtungen der Sonnenwende.}

\begin{center}\fontsize{10}{12}\selectfont\setlength{\tabcolsep}{3pt}
\begin{tabular}{l|c|c|c}
1817 & Anzahl & \shortstack{Wahre Zenitdistanz\\im Meridian} & Im Solstitium \\
\hline
Juni\ 13 & 10 & $28^\circ\ 18'\ 11''30$ & $28^\circ\ 3'\ 52''26$ \\
\ \ \ \ \ \ 16 & 16 & $9\ \ 51.80$ & $48.63$ \\
\ \ \ \ \ \ 17 & 20 & $7\ \ 57.26$ & $55.30$ \\
\ \ \ \ \ \ 18 & 16 & $6\ \ 19.03$ & $53.46$ \\
\ \ \ \ \ \ 19 & 12 & $5\ \ \ 8.69$ & $54.66$ \\
\ \ \ \ \ \ 20 & 20 & $4\ \ 18.74$ & $51.41$ \\
\ \ \ \ \ \ 21 & 12 & $3\ \ 54.41$ & $48.91$ \\
\ \ \ \ \ \ 22 & 10 & $4\ \ \ 0.64$ & $52.10$ \\
\ \ \ \ \ \ 23 & 12 & $4\ \ 28.17$ & $51.77$ \\
\ \ \ \ \ \ 24 & 16 & $5\ \ 20.62$ & $51.62$ \\
\ \ \ \ \ \ 25 & 10 & $6\ \ 36.39$ & $49.96$ \\
\ \ \ \ \ \ 26 & 14 & $8\ \ 20.16$ & $51.61$ \\
\ \ \ \ \ \ 27 & 22 & $10\ 28.75$ & $53.49$ \\
\ \ \ \ \ \ 29 & 16 & $15\ 54.34$ & $51.86$ \\
\ \ \ \ \ \ 30 & 8 & $19\ 16.59$ & $53.88$ \\
\multicolumn{3}{l}{Mittel aus 214 Beobachtungen} & $28^\circ\ 3'\ 52''15$ \\
\end{tabular}
\end{center}

\par\indent Die Zenitdistanz in der ersten Spalte ist lediglich auf den Meridian reduziert und um die Refraktion bereinigt. In der letzten Spalte sind Parallaxe und Einfluss der Sonnenbreite herausgerechnet und die Werte auf den Augenblick der Sonnenwende reduziert. Für die Mitte der Sternwarte erhalte ich also $28^\circ\ 3'\ 52''31$. Mit der Schiefe der Ekliptik $23^\circ\ 27'\ 52''9$ aus den Mailänder Ephemeriden ergibt sich daraus die Polhöhe $51^\circ\ 31'\ 45''21$, um $4''1$ kleiner als aus dem Polarstern. Legt man umgekehrt die Polhöhe aus dem Polarstern zugrunde, so ergibt sich die Schiefe um $4''1$ größer als nach \textsc{Carlini}. Nach \textsc{Bessel} wäre der Unterschied noch anderthalb Sekunden größer. Nach allem, was bisher über dieses \SourceLeaf{407}bemerkenswerte Phänomen gesagt worden ist, scheint es, als wären wir seiner Erklärung nicht näher gekommen. Es wäre gewiss wünschenswert, dass ein geübter Beobachter einmal in einer nicht zu geringen Breite auf der \emph{südlichen} Hemisphäre ähnliche Beobachtungen der Sonne und südlicher Zirkumpolarsterne mit \textsc{Reichenbach}schen Repetitionskreisen anstellt. Für \emph{Reisen} in ferne Weltgegenden zu astronomischen Zwecken, die uns früher so viele wichtige Erfahrungen und Aufschlüsse gebracht haben, haben die Regierungen tatsächlich seit geraumer Zeit gar nichts getan, so sehr sie sonst auch in Europa die Wissenschaft unterstützen.
